% haranoaji-yoko.tex -- pdftex_fontset=haranoaji 横組検査（追加fontset）。
% 判定は runner 側: rc=0、ログマーカ、PDF埋込フォント（pdffonts）が
% HaranoAjiMincho/Gothic（＋実Bold面）、抽出照合。
% 注: \section 見出しを含めると見出し後の本文までゴシックで組まれる
% 既知の縦横backend欠陥（ipa fontsetでも同一; \CJK@global=\global化で
% \gtfamily の family 選択が group を漏れる）があるため、本検査では
% 見出しを使わず明朝本文そのものを検査する。
\RequirePackage{ifpdf}
\documentclass[pdflatex,pdftex_fontset=haranoaji]{jlreq}
\begin{document}
これは原の味フォントセットの横組検証文書である。漢字「括弧」、句読点。
ーっ‥のような縦字形対象文字も横組では通常字形で組まれる。
English混植も含む。あいうえおかきくけこさしすせそたちつてと。
明朝本文の次に太字\textbf{実Bold面}、ゴシック{\gtfamily ゴシック体}、
ゴシック太字{\gtfamily\bfseries ゴシック太字}を検査する。
\typeout{HARANOAJI-YOKO-PASS}
\end{document}
